\icode{N21L-MG-SI-L3-261021-009}
\nceeID{2021山东}{5}
从“玉兔”登月到“祝融”探火，我国星际探测事业实现了由地月系到行星际的跨越。已知火星质量约为月球的9倍，半径约为月球的2倍，“祝融”火星车的质量约为“玉兔”月球车的2倍。在着陆前，“祝融”和“玉兔”都会经历一个由着陆平台支撑的悬停过程。悬停时，“祝融”与“玉兔”所受着陆平台的作用力大小之比为\paren[]
\begin{choices}
	\item \(9:1\)
	\item \(9:2\)
	\item \(36:1\)
	\item \(72:1\)
\end{choices}
\begin{answer}
	B\qquad 提示：黄金代换推出\(g_{mars}:g_{moon} = 9:4\)，并注意到悬停时受力平衡
\end{answer}
